% esl-crossbar-kcl — Analog in-memory-compute crossbar readout: a 2x2 weight-cell array driven by row voltage sources, with per-column KCL current summation feeding a real circuitikz device readout stage.
% THESIS: Many analog per-cell currents reduce to ONE physical KCL summation node per column before any readout device sees them — the column, not the cell, is the unit of analog accumulation.
% Concept grafted from: figstyle.tex conventions (analog-ai EmbML-HB-2026-Ch-AnalogEDAI, grafted, not forked)
% Intent record: intent.yaml (contract §4 shape; ADR-0005 D3 §4 row).
\documentclass[tikz, border=8pt]{standalone}

% --- packages (mirror these in template.meta.json "requires") ---
% circuitikz is the load-bearing addition this template tests: opentikz ships
% ZERO circuitikz templates today (generic TikZ only). See the cp-4856
% learnings note for whether this is a clean graft or fights the model.
\usepackage{circuitikz}
\usepackage{amsmath}
\usetikzlibrary{arrows.meta, backgrounds, positioning, calc}

% --- palette: ESL domain-semantic tokens (contract §2, ADR-0005 D7 -- extends
% opentikz's ot* palette, does not remap onto it; see
% reference/color-palettes/color-palettes.md "ESL domain-semantic tokens").
% Originally grafted from figstyle.tex (analog-ai EmbML-HB-2026-Ch-AnalogEDAI/fig/,
% Gunar-reviewed 2026-07-13 on adcbuffer.tex) as raw blue/green/gray/orange; now
% bound through the named palette tokens tools/validate.py enforces.
\colorlet{inputdomain}{blue}       % inputs / activations
\colorlet{weightdomain}{green}     % analog weight cells
\colorlet{digitaldomain}{gray}     % digital words / logic
\colorlet{analogdomain}{orange}    % analog current / accumulation

\begin{document}
% ==== parameters (edit these) ============================================
\def\ncols{2}          % number of crossbar columns rendered (informational;
                        % this PoC hand-instantiates 2 — see edit_contract
                        % "add-column" for how a real edit extends this)
\def\rowgap{1.1}        % vertical gap between input rows (cm)
\def\colgap{1.6}        % horizontal gap between columns (cm)
% labels
\def\vinonelabel{$V_{1}$}
\def\vintwolabel{$V_{2}$}
\def\kcllabel{$\Sigma$}
\def\outcollabel{$I_{\text{out}}$}
% =========================================================================

\begin{tikzpicture}[
    font=\small,
    >=Stealth,
    % --- grafted node/wire styles, verbatim names from figstyle.tex ---
    kclnode/.style={draw, circle, thick, fill=analogdomain!35,
                    minimum size=6mm, inner sep=0pt},
    cell/.style={draw, rectangle, thick, fill=weightdomain!15,
                 minimum width=0.95cm, minimum height=0.60cm,
                 font=\footnotesize},
    vsource/.style={draw=inputdomain!55, circle, line width=1pt, fill=inputdomain!18,
                    minimum size=8mm, inner sep=0pt,
                    font=\scriptsize\itshape},
    sidenote/.style={font=\fontsize{6.5}{8}\selectfont, digitaldomain!55!black},
    domainlbl/.style={font=\fontsize{7}{9}\selectfont\itshape},
    currentlbl/.style={font=\fontsize{6.5}{8}\selectfont, analogdomain!65!black},
    voltagelbl/.style={font=\fontsize{6.5}{8}\selectfont, inputdomain!65!black},
    awire/.style={analogdomain!80!black, thick},
    dwire/.style={digitaldomain!65!black, thick},
    vwire/.style={inputdomain!60!black, thick},
  ]

  % ---- row voltage sources (analog input activations) ----
  \node[vsource] (v1) at (0, 0)          {\vinonelabel};
  \node[vsource] (v2) at (0, -\rowgap)   {\vintwolabel};

  % ---- crossbar weight cells: 2 rows x 2 cols ----
  \node[cell] (g11) at (\colgap, 0)          {$G_{1,1}$};
  \node[cell] (g21) at (\colgap, -\rowgap)   {$G_{2,1}$};
  \node[cell] (g12) at (2*\colgap, 0)        {$G_{1,2}$};
  \node[cell] (g22) at (2*\colgap, -\rowgap) {$G_{2,2}$};

  \draw[vwire] (v1) -- (g11);
  \draw[vwire] (v1) -- (g12);
  \draw[vwire] (v2) -- (g21);
  \draw[vwire] (v2) -- (g22);

  % ---- per-column KCL summation nodes ----
  \node[kclnode] (k1) at (\colgap, -2*\rowgap-0.4)   {\kcllabel};
  \node[kclnode] (k2) at (2*\colgap, -2*\rowgap-0.4) {\kcllabel};

  \draw[awire] (g11) -- (k1);
  \draw[awire] (g21) -- (k1);
  \draw[awire] (g12) -- (k2);
  \draw[awire] (g22) -- (k2);

  \node[currentlbl, left=1pt of k1.west, xshift=-2pt] {$I_{1,1}{+}I_{2,1}$};
  \node[currentlbl, right=1pt of k2.east, xshift=2pt] {$I_{1,2}{+}I_{2,2}$};

  % ---- real circuitikz device symbols on the readout path (the genuine
  %      EE-symbol test: opentikz's generic-TikZ node styles above cannot
  %      express these; only circuitikz's `to[...]` bipole syntax can) ----
  \coordinate (k1out) at ($(k1)+(0,-1.15)$);
  \coordinate (k2out) at ($(k2)+(0,-1.15)$);
  \draw (k1) to[R=$R_{\text{sense}}$, o-o] (k1out);
  \draw (k2) to[I, l=$I_{\text{ref}}$, o-o] (k2out);

  \node[sidenote, below=2pt of k1out, xshift=-9pt] {col.\,1 readout};
  \node[sidenote, below=2pt of k2out, xshift=9pt] {col.\,2 readout};

  \node[domainlbl, inputdomain!65!black, above=6pt of v1] {Analog domain};
  \node[currentlbl, right=4pt of k2out] {\outcollabel};

\end{tikzpicture}
\end{document}
